\section{Three smaller Jacobians for the actual first-square locus}
\label{hg-section}

Fix an odd integer $g\geq3$ and $u\in\{1,\zeta,\zeta^2\}$.
Write $A=A_{E,u}(S,T)$ and $B=B_{E,u}(S,T)$ for the homogeneous
degree-$g$ integer forms in Section~\ref{rl-section}, so that
$A+B\zeta=u(S+T\zeta)^g$, and put $r=A/B$.
The curve $D=D_{g,u}$ is the smooth projective normalization of
\begin{equation}\label{hg-original}
 r(s)=\frac{a_0(t)b_0(t)}{U_0(t)^2},\qquad
 a_0=t^2-1,\quad b_0=2t+1,\quad U_0=t^2+t+1.
\end{equation}

\begin{theorem}[Three hyperelliptic quotients]\label{hg-three-quotients}
There is an equality of function fields over $\mathbb Q$
\begin{equation}\label{hg-biquadratic}
 \mathbb Q(D)=\mathbb Q(s)(v,w),\qquad
 v^2=1+r,\quad w^2=1-3r.
\end{equation}
This is a geometrically connected $V_4$ cover of $\mathbb P^1_s$,
of degree four and genus $3g-3$. Its three quadratic quotient
curves have homogeneous hyperelliptic equations
\begin{equation}\label{hg-hyperelliptic}
 H_1:\ y_1^2=B(A+B),\qquad
 H_2:\ y_2^2=B(B-3A),\qquad
 H_3:\ y_3^2=(A+B)(B-3A).
\end{equation}
Each right side is a squarefree binary form of degree $2g$;
the $y_j$ have weight $g$, and the equations denote their smooth
projective models. Each $H_j$ has genus $g-1$. On a dense affine
chart the quotient maps are
\[
 y_1=Bv,\qquad y_2=Bw,\qquad y_3=Bvw.
\]
\end{theorem}
\begin{proof}
Put $c_0=t^2+2t$ and $d_0=t^2-2t-2=a_0-b_0$.
The identities
\[
 c_0^2=U_0^2+a_0b_0,\qquad d_0^2=U_0^2-3a_0b_0
\]
give $v=c_0/U_0$ and $w=d_0/U_0$. Conversely the inverse on a
dense open is
\begin{equation}\label{hg-inverse}
 t=\frac{v+w+2}{v-w}.
\end{equation}
Indeed $w^2+3v^2=4$ and
$a_*=(v+w)/2$, $b_*=(v-w)/2$ satisfy
$a_*^2+a_*b_*+b_*^2=1$. Substitution gives
$a_0/b_0=a_*/b_*$ and $U_0/b_0=1/b_*$, recovering both $v,w$
and $r=a_*b_*$. This proves the function-field identity, which
extends to the smooth projective models
\cite{StacksCurves2026}. The dense-open formulas do not discard
exceptional points from that projective assertion.

Over $K_E$, the map $r(s)$ is conjugate by $\rho_E$ to a nonzero
scalar times a $g$-th power. Its only branch values are
$-\zeta,-\bar\zeta$. The rational target values
$\infty,-1,1/3$ are distinct and avoid those values. Their
inverse images are therefore three disjoint sets of $g$ simple
points. They are exactly the zeros of $B,A+B,B-3A$, including
any zero at infinity in homogeneous coordinates.

At a zero of $A+B$ the valuation of $1+r$ detects its nontrivial
square class, and at a zero of $B-3A$ the valuation of $1-3r$
detects the other class. They are independent, so the extension
in \eqref{hg-biquadratic} has geometric group $V_4$.
At zeros of $A+B$ only $v$ ramifies, at zeros of $B-3A$ only
$w$ ramifies, and at zeros of $B$ their common simple poles
give the diagonal inertia subgroup of order two. There are no
other branch points. Thus
\[
 2g(D)-2=-8+(3g)\,2=6g-8.
\]
The three nontrivial square classes give the three quadratic
subfields. Clearing by the square $B^2$ gives
\eqref{hg-hyperelliptic}. Each has $2g$ simple branch points,
so $2g(H_j)-2=-4+2g$.
This argument includes infinity and the normalization at common
branch points of the raw fiber product of $H_1$ and $H_2$.
\end{proof}

\begin{theorem}[An isogeny over $\mathbb Q$ and transport of odd torsion]
\label{hg-isogeny}
Let $p_j:D\to H_j$ be the quotient maps, with Jacobians $J_D,J_j$.
Define homomorphisms over $\mathbb Q$ by
\begin{align}
 \Phi:J_D&\longrightarrow J_1\times J_2\times J_3,
 &P&\longmapsto ((p_1)_*P,(p_2)_*P,(p_3)_*P),\label{hg-norm-map}\\
 \Psi:J_1\times J_2\times J_3&\longrightarrow J_D,
 &(P_1,P_2,P_3)&\longmapsto
 p_1^*P_1+p_2^*P_2+p_3^*P_3.\label{hg-pullback-map}
\end{align}
Then $\Psi\Phi=[2]$, and $\Psi$ is an isogeny over $\mathbb Q$.
For $u=1$, the class $\Theta$ of
Theorem~\ref{ug-exact-torsion} maps under $\Phi$ to a
$K$-rational point of exact order $g$ on the product.
For prime $g$, at least one coordinate has exact order $g$.
\end{theorem}
\begin{proof}
Finite maps of smooth curves are finite locally free. Pullback
and norm of line bundles therefore induce the displayed
homomorphisms on Jacobians. On divisors these are the usual
pullback and pushforward, with norms of rational functions
ensuring that principal divisors map to principal divisors.
The constructions and all $p_j$ are defined over $\mathbb Q$;
this is the standard Picard/Jacobian functoriality
\cite{MilneJV2021}.

If $\sigma_j$ generates the group of $D\to H_j$, the geometric
divisor identity is $p_j^*(p_j)_*=1+\sigma_j$.
At ramified points it still holds: the pullback multiplicity
is two and the two right-hand terms coincide.
For $q:D\to\mathbb P^1_s$, the analogous identity is
\[
 q^*q_*=1+\sigma_1+\sigma_2+\sigma_3.
\]
It vanishes on the Jacobian because
$\operatorname{Pic}^0(\mathbb P^1)=0$. Hence
\begin{equation}\label{hg-times-two}
 \Psi\Phi=3+\sigma_1+\sigma_2+\sigma_3=[2].
\end{equation}
The equality of homomorphisms holds over $\mathbb Q$ since
it holds after algebraic closure. Multiplication by two is
surjective, so $\Psi$ is surjective. The source and target
dimensions are both $3(g-1)$, and its kernel is therefore
finite. This proves the isogeny.

The point $\Phi(\Theta)$ is killed by $g$. If it is killed by
an integer $n$, \eqref{hg-times-two} shows that $2n$ kills
$\Theta$. Since $\Theta$ has exact odd order $g$, this forces
$g\mid n$. Thus the tuple has exact order $g$. For prime $g$
at least one nonzero coordinate has that order; for composite
$g$ the tuple's exact order need not occur in one coordinate,
and no such stronger statement is made.
\end{proof}

\begin{theorem}[The positive rational chart]\label{hg-positive}
On the nonbranch rational chart, the points of $D$ with $t>1$
are exactly the rational triples $(s,v,w)$ in
\eqref{hg-biquadratic} satisfying
\begin{equation}\label{hg-real-gate}
 0<r(s)<1/3,\qquad v>0.
\end{equation}
\end{theorem}
\begin{proof}
These conditions imply $1<v<2/\sqrt3$ and $|w|<1$.
The denominator $v-w$ in \eqref{hg-inverse} is positive,
and the numerator exceeds it by $2(w+1)>0$, so $t>1$.
Conversely rational $t>1$ gives $v=c_0/U_0>0$ and
$r=a_0b_0/U_0^2>0$. The identity for $d_0^2$ gives
$r\leq1/3$. Equality would require
$t^2-2t-2=0$, whose roots $1\pm\sqrt3$ are irrational.
Thus the inequality is strict. The dense-open inverse formulas
are defined throughout this chart, and keep the ordered sign
choice of $w$.
\end{proof}

By Section~\ref{rl-section}, these points give actual primitive
seeds with square first norm and $g$-th-power second norm.
When $\gcd(g,6)=1$, the single $u=1$ curve gives the
ordered-seed bijection. For $3\mid g$, retain all three unit
choices. A point on one quotient alone is merely necessary:
the points on $H_1,H_2$ must have the same source coordinate
$s$ and satisfy \eqref{hg-real-gate}.

For the explicit case $g=3,u=1$, the affine forms are
\begin{gather}
 A=s^3-3s-1,\qquad B=3s(s+1),\label{hg-cubic-forms}\\
 A+B=s^3+3s^2-1,\qquad
 B-3A=-3s^3+3s^2+12s+3.\nonumber
\end{gather}
They define three genus-two curves in
\eqref{hg-hyperelliptic}, whose Jacobian product is isogenous
to the genus-six Jacobian of $D_{3,1}$. The homogeneous
degree-three $B$ includes a simple zero at infinity.
The other unit choices follow by integer coordinate
multiplication and have not been discarded.
This explicit reduction supplies a smaller computation problem.
It is not a rank computation or a determination of all rational
points on any quotient. No fixed-exponent solution, uniform
height bound, or control of varying residuals follows without
additional work.
